\section{DeepSeekMath-V2：让 verifier 也接受验证}

\subsection{Outcome reward 无法保证证明过程}

\videofigure{figures/outcome-reward-models.jpg}{Outcome reward 只检查最终答案，无法判断中间证明是否存在跳步、循环论证或幸运猜中。}{00:15:02--00:16:52}

\videofigure{figures/proof-verification-problem.jpg}{LLM 可生成形式流畅但数学无效的证明，而没有参考解时，普通 judge 也难稳定定位问题。}{00:16:52--00:18:18}

数学证明不是短答案匹配。一个 theorem proof 可能结论正确，但某一步不从前提推出；只给终局 reward 会把这种 flawed trajectory 当作正样本。Process reward model 又依赖昂贵的逐步人工标签。

\begin{warningbox}{Verifier 可能“分数对、理由错”}
LLM judge 即使给出大致正确的总分，也可能编造不存在的错误。若这些 critique 被用于训练 generator，系统会把 verifier 的幻觉写回模型。
\end{warningbox}

\subsection{Generator--Verifier--Meta-Verifier}

\videofigure{figures/deepseekmath-v2-architecture.jpg}{DeepSeekMath-V2 构建 generator、verifier 与 meta-verifier 的协同闭环。}{00:18:18--00:19:34}

结构可以形式化为：

$$
y\sim G_\theta(x),\qquad
(e,s)=V_\phi(x,y),\qquad
m=M_\psi(x,y,e,s),
$$

其中 generator 产生证明 $y$，verifier 给出问题分析 $e$ 和分数 $s$，meta-verifier 判断“问题是否真实存在、分数是否由分析推出”。

\videofigure{figures/meta-verification.jpg}{Meta-verifier 对 verifier 的分析做二次检查，降低 fabricated issue，并提升解释可信度。}{00:19:34--00:20:18}

\begin{importantbox}{Meta-verification 检查的是评价的证据链}
它不只是再投一次票，而是验证 critique 中每个 issue 是否能在原证明中定位，以及最终 score 是否与这些 issue 一致。
\end{importantbox}

\subsection{迭代优化与结果}

\videofigure{figures/evaluation-methodology.jpg}{系统结合 GRPO、迭代生成--验证，以及多类数学证明 benchmark 评估自验证能力。}{00:20:18--00:20:52}

\videofigure{figures/proof-results.jpg}{在 IMO Shortlist 2024 上，pass@1 与 best@32 都随多轮 self-verification 迭代持续上升。}{00:20:52--00:21:22}

\videofigure{figures/self-verification-improvement.jpg}{自选高 verifier score 的证明用于下一轮训练，形成 generator 与 verifier 共同爬坡。}{00:21:22--00:22:22}

若从 $K$ 个证明中选择 verifier 分数最高者：

$$
y^*=\arg\max_{y_k,\,k\le K}V_\phi(x,y_k),
$$

则前提是 $V_\phi$ 的 ranking 与真实证明质量一致。Meta-verification 的作用正是改善这一 alignment，使 best@$K$ 不只是挑出“最会讨好 judge”的证明。

\begin{warningbox}{更多 judge 层不自动带来真理}
Generator、verifier 和 meta-verifier 可能共享训练数据与偏见，产生 correlated error。仍需要形式化检查、独立模型、专家审计或 adversarial test 来打破共谋。
\end{warningbox}

\subsection{本章小结}

DeepSeekMath-V2 把 verifier 从不可质疑的裁判改造成可审计对象。它展示了没有参考证明时的自动过程验证路径，但仍依赖数学领域相对清晰的逻辑结构。

